\chapter{Preliminaries}
\label{ch:preliminaries}
% Definitions and notation shared by all chapters, known results with citations; the list of symbols points here.
% Guidance: knowledge/writing/thesis.md section 4, knowledge/writing/math-writing.md

Preliminaries text.
Graph-theoretic terminology follows~\cite{Diestel2025}; algebraic methods are surveyed in~\cite{GodsilRoyle2001}.

% --- Example (remove): definition, figure, table --------------------------------------------
% Labels: prefix + lowercase words with hyphens (ch:, sec:, def:, fig:, tab:); \label directly after \caption.
%
% \begin{definition}[Cut-path]
% \label{def:cut-path}
% Let $G = (V, E)$ be a graph and $u, v \in V$.
% A set $S \subseteq E$ is a \emph{cut-path} if \dots
% \end{definition}
%
% \begin{figure}[tb]
%   \centering
%   \includegraphics[width=0.6\linewidth]{img/example}   % vector PDF; file name lowercase-with-hyphens
%   \caption{What the reader should see, in one sentence.}
%   \label{fig:example}
% \end{figure}
%
% \begin{table}[tb]
%   \centering
%   \caption{Caption above the table. Running time in seconds, mean of 10 runs.}
%   \label{tab:runtime}
%   \begin{tabular}{@{}lrr@{}}
%     \toprule
%     Instance & $\abs{V}$ & Time (s) \\
%     \midrule
%     \texttt{grid-10} & 100 & 0.12 \\
%     \texttt{grid-20} & 400 & 1.85 \\
%     \bottomrule
%   \end{tabular}
% \end{table}
%
% \Cref{def:cut-path} fixes the notation; \cref{fig:example} shows \dots; \cref{tab:runtime} lists \dots
